\chapter{서론}
\section{연구 배경}
본 과제의 목표는 SO-ARM101 로봇팔로 다섯 블록을 제한시간 안에 지정 영역으로 옮기고, 후속 미션에서 적재하는 것이다. 이하에서는 SO-ARM101을 SO-101로 표기한다. 초기에는 이동형 로봇의 물체 수집을 검토했으나 과제 규격이 바뀌면서 SO-101 블록 조작으로 범위를 변경하였다. 개발 환경에서 개별 동작을 성공시키는 것과 함께, 장비를 실제 시연장으로 옮긴 뒤에도 여러 블록을 반복 처리할 수 있어야 했다.

초기에는 리더암으로 단일 블록을 집어 중앙 목표 지점으로 옮기는 동작을 시연하고, 이를 ACT와 SmolVLA로 학습하였다. 시작 위치를 한정하면 수행이 나아졌지만 작업 영역 전체로 범위를 넓혔을 때는 안정적인 동작을 얻기 어려웠다. ACT와 컴퓨터 비전·역기구학(CV+IK)을 결합한 방식에서도 카메라 위치 변화에 따른 성능 변동과 시연 재수집 부담이 남았다.

이에 크기와 색상이 정해진 블록과 고정 작업 평면의 조건을 활용해 CV+IK 기반 제어를 구성하였다. 검출 위치, 좌표 변환, 파지 확인과 재시도를 나누어 점검하면서 동작을 보완하였으며, 반복 실행을 학습 자료로 저장하여 수동 시연 수집의 부담도 줄이고자 하였다.

\section{초기 접근의 한계와 설계 요구}
첫째, 특정 시작 위치에서 얻은 결과를 작업 영역 전체의 수행으로 확장할 수 있어야 했다. 단일 블록이라도 위치와 접근 조건이 달라질 수 있으므로, 검출된 위치에서 실행 가능한 파지 자세를 생성해야 한다.

둘째, 환경 변화의 영향을 확인하고 조정할 수 있어야 했다. 카메라 고정·재캘리브레이션과 조명에 따른 검출 조정은 CV+IK에서도 필요하다. 검출 중심, 로봇 좌표, 목표 자세와 센서 판정을 각각 확인하여 수정할 단계를 좁히도록 하였다.

셋째, 개별 파지 실패를 전체 미션의 중단으로 이어지지 않게 해야 했다. 제한시간 내 여러 블록을 처리하려면 파지를 확인한 뒤 운반하고, 실패 시 재시도하거나 대상을 다시 선택해야 한다. 최종 제어 흐름은 이러한 판단과 복구 절차를 명시적으로 관리하도록 구성하였다.

\section{연구 목표}
\subsection{미션 규격과 성공 조건}
블록은 $40\times40\times20\,\mathrm{mm}$ 크기이며, 서로 다른 색의 블록 다섯 개를 사용한다. 미션 시작 시 각 블록은 넓은 면을 바닥에 둔 채 지정 영역 밖에 놓인다. 블록끼리 맞닿을 수는 있지만 위아래로 겹치지는 않는다. 두 미션은 블록을 다시 배치한 뒤 각각 수행한다. 목표와 평가 조건은 표~\ref{tab:mission}과 같다.

\begin{table}[htbp]
\centering
\caption{과제 미션의 목표와 평가 조건}\label{tab:mission}
\begin{tabularx}{\linewidth}{@{}L{0.16\linewidth}YY@{}}
\toprule
항목 & 1차 미션: 이동 & 2차 미션: 적재 \\
\midrule
목표 & 다섯 개 블록을 지정 영역에 배치 & 다섯 개 블록을 수직으로 적재 \\
제한시간 & 180초 & 300초 \\
지정 영역 & \multicolumn{2}{l}{$200\times100\,\mathrm{mm}$ 내부 영역} \\
인정 조건 & 경계 20\,mm 이내 걸침, 세로 세움과 일부 겹침 허용. 적재는 불인정 & 적재 상태를 5초 이상 안정적으로 유지 \\
부분 수행 & 시간 종료 시 영역에 놓인 블록 수 & 시간 종료 시 적재한 블록 수 \\
\bottomrule
\end{tabularx}
\end{table}

지정 영역은 로봇 베이스에서 약 25\,cm 떨어진 곳에 두며, 제어할 때는 네 꼭짓점을 로봇 베이스 좌표계에 등록한다.

신경망은 최대 하나까지 사용할 수 있으나 필수는 아니다. 최종 CV+IK 시스템은 신경망 없이 영상 처리와 기구학 계산으로 동작한다.

\subsection{세부 연구 목표}
영상에서 지정 영역 밖의 블록을 골라 로봇 좌표로 변환하고, 관절 한계와 그리퍼 방향을 고려한 접근·파지·이동 자세를 생성한다. 파지 확인과 재시도를 포함한 상태 흐름으로 블록을 반복 처리하며, 단계별 자료를 기록해 실패 지점을 추적한다. 적재 미션으로 동작을 확장하고, CV+IK의 실행 자료를 학습 데이터로 저장하여 향후 모방학습을 다시 시도할 기반을 마련한다.

본 연구는 기존 영상 처리와 기구학 도구를 과제 조건에 맞게 통합하고, 실패 복구와 실행 자료 수집까지 연결하는 데 초점을 두었다. 2장은 관련 기술, 3장은 설계와 구현, 4장은 실험 결과와 한계를 설명하며, 5장에서 결론과 후속 연구를 제시한다.
\FloatBarrier
